% !TITLE 归园田居·其一
% !AUTHOR 陶渊明
% !DYNASTY 魏晋
% !GENRE 五言古诗
% !ORDER 1

\begin{poem}
少无适俗韵，性本爱丘山。
误落尘网中，一去三十年。
羁鸟恋旧林，池鱼思故渊。
开荒南野际，守拙归园田。
方宅十余亩，草屋八九间。
榆柳荫后檐，桃李罗堂前。
暧暧远人村，依依墟里烟。
狗吠深巷中，鸡鸣桑树颠。
户庭无尘杂，虚室有余闲。
久在樊笼里，复得返自然。
\end{poem}

\begin{appreciation}
这是陶渊明最著名的诗作之一，写他辞去彭泽县令、回归田园时的心情。

陶渊明生活在东晋末年，那是一个政治黑暗、战乱频仍的时代。他做过几任小官，每次都不长久。最后一次是彭泽县令，上任八十多天，上级派督邮来检查，县吏说应当束带迎接，陶渊明叹道：我岂能为五斗米折腰向乡里小儿！当天就解印绶辞职，从此再不出仕。

少无适俗韵，性本爱丘山——从小就没有适应世俗的气质，天性热爱山林。误落尘网中，一去三十年——错误地落入尘世的罗网，一去就是三十年。尘网和樊笼是陶渊明对官场的比喻，可见他对仕途的厌恶是发自内心的。

羁鸟恋旧林，池鱼思故渊——被束缚的鸟思念旧林，池中的鱼想念深潭。这是全诗最精彩的比喻。鸟和鱼本属于自然，却被困在人为的环境中；人也是一样，本性热爱自由，却被官场的规矩束缚。这个比喻既形象又贴切，成为后世表达归隐之情的经典意象。

方宅十余亩，草屋八九间以下八句，写田园生活的宁静美好。暧暧远人村，依依墟里烟——远处的村庄隐约可见，炊烟轻柔地上升。暧暧（ài ài）是模糊的样子，依依是轻柔的样子。这两句没有主语，只有画面，像一幅水墨山水。狗吠深巷中，鸡鸣桑树颠——深巷里传来狗叫声，桑树顶上公鸡在打鸣。这是田园生活的声音，充满了生机和烟火气。

最后两句久在樊笼里，复得返自然是全诗的总结。从尘网到樊笼，从误落到复得，陶渊明完成了一个从迷失到回归的精神旅程。他的归隐不是消极的逃避，而是一种主动的选择——在乱世中守护内心的自由。这种选择，让后世的无数文人在困顿中看到了另一种可能性。
\end{appreciation}
